% Preámbulo del informe: paquetes, colores, títulos, código y figuras.

\usepackage[T1]{fontenc}
\usepackage[spanish,es-noquoting,es-tabla]{babel}
\usepackage[a4paper,margin=2.1cm,top=2.3cm,bottom=2.2cm]{geometry}
\usepackage{amsmath,amssymb}
\usepackage{booktabs}
\usepackage{array}
\usepackage{xcolor}
\usepackage{tikz}
\usepackage{float}
\usepackage{listings}
\usepackage{tcolorbox}
\usepackage{enumitem}
\usepackage{fancyhdr}
\usepackage{titlesec}
\usepackage[font=small,labelfont=bf,skip=4pt]{caption}
\usepackage[hidelinks]{hyperref}

\usetikzlibrary{positioning,arrows.meta}
\tcbuselibrary{skins,breakable}

% --- paleta ---------------------------------------------------------------
\definecolor{tinta}{HTML}{1B2733}
\definecolor{acento}{HTML}{1F6FEB}
\definecolor{acentoSuave}{HTML}{DDE9FB}
\definecolor{exito}{HTML}{1A7F37}
\definecolor{exitoSuave}{HTML}{DCF4E3}
\definecolor{gris}{HTML}{6E7781}
\definecolor{grisSuave}{HTML}{F0F2F5}

\color{tinta}

% --- títulos, párrafos y encabezados ---------------------------------------
\titleformat{\section}{\Large\bfseries\color{acento}}{\thesection}{0.7em}{}
\titleformat{\subsection}{\large\bfseries}{\thesubsection}{0.6em}{}
\titlespacing*{\section}{0pt}{14pt}{6pt}
\titlespacing*{\subsection}{0pt}{10pt}{4pt}
\titlespacing*{\paragraph}{0pt}{7pt}{0.8em}
\setlength{\parskip}{4pt}
\setlength{\parindent}{0pt}

\pagestyle{fancy}
\fancyhf{}
\fancyhead[L]{\small\color{gris}Minigestor de base de datos multimodal}
\fancyhead[R]{\small\color{gris}Informe: Partes 1 y 2}
\fancyfoot[C]{\small\color{gris}\thepage}
\renewcommand{\headrulewidth}{0.4pt}

% --- tablas ----------------------------------------------------------------
% Columna de párrafo sin justificar: en celdas estrechas la justificación parte palabras.
\newcolumntype{L}[1]{>{\raggedright\arraybackslash}p{#1}}

% --- código ---------------------------------------------------------------
\lstdefinestyle{consola}{
  basicstyle=\ttfamily\footnotesize,
  keywordstyle=\color{acento}\bfseries,
  commentstyle=\color{gris}\itshape,
  stringstyle=\color{exito},
  numbers=none,
  frame=leftline,
  framesep=8pt,
  rulecolor=\color{acentoSuave},
  framerule=2pt,
  backgroundcolor=\color{grisSuave},
  xleftmargin=10pt,
  breaklines=true,
  showstringspaces=false,
  columns=fullflexible,
  keepspaces=true
}
\lstdefinestyle{sql}{
  style=consola,
  morekeywords={SELECT,FROM,WHERE,CREATE,TABLE,INDEX,USING,ORDER,BY,LIMIT,ON,AND,RTREE}
}
\lstset{style=consola}

% --- cajas y figuras --------------------------------------------------------
\newtcolorbox{resumen}{
  enhanced, colback=exitoSuave, colframe=exito,
  boxrule=0pt, leftrule=3pt, arc=2pt, left=8pt, right=8pt, top=5pt, bottom=5pt}

\tikzset{
  flecha/.style={-{Stealth[length=2mm]}, gris, thick},
}

% Una gráfica de benchmarks/figures con su pie: [ancho]{archivo}{etiqueta}{pie}.
\newcommand{\grafica}[4][\textwidth]{%
  \begin{figure}[!htb]
    \centering
    \includegraphics[width=#1]{../../benchmarks/figures/#2}
    \caption{#4}
    \label{#3}
  \end{figure}}
